\hojaejercicio{Convergencia de seires complejas y series de potencias}
%ejercicio1
\ejercicio{Estudia la convergencia y la convergencia absoluta de $\sum_{n=1}^{\infty}\frac{1+2i}{2^{n}}$; $\sum_{n=1}^{\infty}\left(\frac{i}{2}+\frac{i}{2}i\right)^{n}$; $\sum_{n=2}^{\infty}\left(\frac{\log n}{n}+i^{n}\frac{\log n}{n}\right)$.}{
    La serie converge absolutamente.
\[ \sum_{n=1}^{\infty}|\frac{1+2i}{2^{n}}| \text{ converge } \iff \sum_{n=1}^{\infty}(|1+2i|)\left(\frac{1}{2}\right)^{n} \text{ converge } \iff \sum_{n=1}^{\infty}\left(\frac{1}{2}\right)^{n} \text{ converge.} \]
Como $\sum_{n=1}^{\infty}r^{n}$ converge cuando $0 \le r < 1$, entonces la serie dada converge absolutamente y por tanto, converge.

 $\sum_{n=2}^{\infty}\left(\frac{\log n}{n}+i^{n}\frac{\log n}{n}\right)$
Consideramos la serie $\sum_{n=2}^{\infty}(1+i^{n})\frac{\log n}{n} = \sum_{n=2}^{\infty}a_{n}$ siendo $a_{n} = (1+i^{n})\frac{\log n}{n}$. Analizamos la parte real $Re(a_{n})$ según $n$:
\begin{itemize}
    \item Si $n=1+4k$ para algún $k \in \mathbb{Z}$: $Re(a_{n}) = \frac{\log n}{n}$
    \item Si $n=2+4k$ para algún $k \in \mathbb{Z}$: $Re(a_{n}) = 0$
    \item Si $n=3+4k$ para algún $k \in \mathbb{Z}$: $Re(a_{n}) = \frac{\log n}{n}$
    \item Si $n=4k$ para algún $k \in \mathbb{Z}$: $Re(a_{n}) = \frac{2 \log n}{n}$
\end{itemize}
Como $\sum_{k=1}^{\infty}\frac{1}{1+4k}$ (que diverge) $\le \sum_{k=1}^{\infty}\frac{\log(1+4k)}{1+4k}$, la subserie diverge. Esto implica que la parte real diverge, por lo que la serie original diverge.
}

%ejercicio2
\ejercicio{Determina el radio de convergencia de: $\sum_{n=1}^{\infty} n^{p}z^{n}$ ; $\sum_{n=1}^{\infty} q^{n^{2}}z^{n}$; $\sum_{n=1}^{\infty}z^{n!}$ $\sum_{n=1}^{\infty}n!z^{n}$; $\sum_{n=1}^{\infty}\frac{n!}{n^{n}}z^{n}$.}{
    \subsection*{a) $\sum_{n=1}^{\infty} n^{p}z^{n}$}
Usando el criterio de la raíz:
\[ \limsup_{n\to\infty} \sqrt[n]{n^{p}} = 1 \implies R=1 \]

\subsection*{b) $\sum_{n=1}^{\infty} q^{n^{2}}z^{n}$}
Usando el criterio de la raíz:
\[ \limsup_{n\to\infty} \sqrt[n]{q^{n^{2}}} = \limsup_{n\to\infty} q^{n} = 0 \implies R=+\infty \]

\subsection*{e) $\sum_{n=1}^{\infty} \frac{n!}{n^{n}}z^{n}$}
Usando el criterio del cociente:
\[ \lim_{n\to\infty} \frac{\frac{n!}{n^{n}}}{\frac{(n+1)!}{(n+1)^{n+1}}} = \lim_{n\to\infty} \frac{n!(n+1)^{n+1}}{n^{n}(n+1)!} = \lim_{n\to\infty} \left(\frac{n+1}{n}\right)^{n} = e \implies R=e \]
}

%ejercicio3
\ejercicio{Calcula $\sum_{n=1}^{\infty}nz^{n}$ y $\sum_{n=1}^{\infty}n^{p}z^{n}$.}{
    Usando el criterio de la raíz se comprueba que $R=1$. Consideremos $g(z) = \sum_{n=1}^{\infty} z^{n}$.
\[ f(z) = z \cdot g'(z) = z \cdot \left(\sum_{n=1}^{\infty}z^{n}\right)' = z \cdot \left(\frac{z}{1-z}\right)' = z \cdot \left(\frac{1}{(1-z)^{2}}\right) = \frac{z}{(1-z)^{2}} \]
}

%ejercicio4
\ejercicio{Si $R$ es el radio de convergencia de $\sum_{n=0}^{\infty}a_{n}z^{n}$, ¿cuál es el radio de convergencia de $\sum_{n=0}^{\infty}a_{n}^{2}z^{n}$ y de $\sum_{n=0}^{\infty}a_{n}z^{2n}$? }{}

%ejercicio5
\ejercicio{Sea $m \in \mathbb{N}$, expresa $\frac{1}{(1-z)^{m}}$ en serie de potencias de $z$.}{
    Si $f(z) = \frac{1}{(1-z)^{m}} = \sum_{n=0}^{\infty}a_{n}z^{n}$, entonces $f^{(n)}(0) = n! a_{n} \implies a_{n} = \frac{f^{(n)}(0)}{n!}$.
Se comprueba que:
\[ \frac{f^{(n)}(0)}{n!} = \binom{m+n-1}{n} \]
Luego:
\[ g(z) = 1 + \sum_{n=1}^{\infty} \binom{m+n-1}{n} z^{n} \]
}

%ejercicio6
\ejercicio{Expresa $\frac{2z+3}{z+1}$ en serie de potencias de $z-1$. ¿Cuál es el radio de convergencia de esta serie de potencias? }{
    \[ \frac{2z+3}{z+1} = 2 + \frac{1}{z+1} = 2 + \frac{1}{2+(z-1)} = 2 + \frac{1}{2} \cdot \frac{1}{1+\frac{z-1}{2}} \]
\[ = 2 + \frac{1}{2} \sum_{n=0}^{\infty} \frac{(-1)^{n}}{2^{n}}(z-1)^{n} \]
El radio de convergencia es:
\[ \frac{1}{R} = \limsup_{n\to\infty} \sqrt[n]{\frac{1}{2^{n}}} = \frac{1}{2} \implies R=2 \]
}

%ejercicio7
\ejercicio{¿Para qué valores de $z$ converge la serie $\sum_{n=0}^{\infty}\left(\frac{z}{1+z}\right)^{n}$? }{}

%ejercicio8
\ejercicio{Si $f(z)=\sum_{n=0}^{\infty}a_{n}z^{n}$, ¿qué valores toma $\sum_{n=0}^{\infty}na_{n}z^{n}$? }{}

%ejercicio9
\ejercicio{Sea $f(z)=\begin{cases}\frac{e^{z}-1}{z}, & \text{si } z \neq 0 \\ 1, & \text{si } z = 0 \end{cases}$. Halla la representación de $f$ en serie de potencias alrededor de 0 y determina el radio de convergencia de la serie.}{}

%ejercicio10
\ejercicio{Sea $f$ la función suma de la serie de potencias $\sum_{n=0}^{\infty}a_{n}z^{n}$ de radio de convergencia $R > 0$. Si existe $(z_{k})_{k=1}^{\infty} \subset D(0;R) \setminus \{0\}$ convergente a 0 tal que $f(z_{k})=0$ para todo $k \in \mathbb{N}$, demuestra que $f \equiv 0$ (Sugerencia: pruébalo por inducción que $a_{n}=0$ para todo $n \in \mathbb{N}$).}{}